\documentclass[../main]{subfiles}
\begin{document}
\ifSubfilesClassLoaded{\setcounter{page}{62}}{}
\section{Conclusion}
\label{sec:11_conclusion}

Only now, having outlined in broad terms the main elements of communism and shown the seeds from which it can grow into a social system, do we consider ourselves entitled to move on to definitions. Otherwise they would be dead, meaningless abstractions rather than an answer to the question posed in the title: “What is communism?”

Communism (communist society) is a free association of people who rationally create their own lives and consciously posit their own free activity.

Communism is the directly collective form of free activity, having become total as the form of a society that creates its own development.

Communism is, as Marx and Engels put it in the “Manifesto”, “an association, in which the free development of each is the condition for the free development of all”.

The time has come to repeat once again a slogan that has not aged one bit:

WORKERS OF ALL COUNTRIES, UNITE!
\end{document}
